% ============================================================================
%  SECTION 5 — Homogenized constrained mixture
% ============================================================================
\begin{frame}[t]{}
  \vfill\begin{center}{\Huge\color{YaleBlue} Homogenized CMM}
  \end{center}\vfill
\end{frame}

\begin{frame}[t]{Throw Away the History, Keep Its Average}
  \textbf{Temporal homogenization} replaces the entire deposition history by a
  \emph{single evolving mean natural configuration} per constituent --- two ODEs
  per constituent, cost $O(N^2)\to O(N)$.
  \dlab{general form (Cyron 2016) --- mean natural-config tensor}
  \zlab{reduction --- scalar mean stretch}
  \vspace{-0.2em}
  \giveneq{\bm{F}^k_{\text{e}} = \bm{F}(s)\,\bm{F}^{k}_{n}(s)^{-1}
    \qquad\longrightarrow\qquad \lambda_e^k = \frac{\lambda}{\lambda_n^k}.}
  \vspace{-0.2em}
  \begin{center}\homogmap[0.70]\end{center}
  \srccite{Cyron, Aydin, Humphrey, Biomech Model Mechanobiol (2016)}
\end{frame}

% ---- Maxwell analog moved EARLY: the visual intuition for the whole model ----
\begin{frame}[t]{The Mechanical Analog: A Maxwell Fluid with a Motor}
  Each constituent behaves like a \textbf{Maxwell element in parallel with a
  motor}: turnover relaxes stress toward a \emph{prestress}, not zero.
  \begin{center}\maxwellelement[0.78]\end{center}
  \begin{columns}[T]
    \begin{column}{0.52\textwidth}
      \footnotesize
      \renewcommand{\arraystretch}{1.5}
      \begin{tabular}{@{}ll@{}}
        $\sigma^i$ & constituent stress (the load)\\[2pt]
        $\mathbb{C}^{i}$ & elastic stiffness (spring)\\[2pt]
        $\dot\varrho_0^{\,i}/\varrho_0^{\,i}+1/T^i$ & turnover + degradation (dashpot)\\[2pt]
        $\bar\sigma^{i}_{\text{pre}}$ & deposition prestress (motor)\\
      \end{tabular}
    \end{column}
    \begin{column}{0.48\textwidth}
      \footnotesize \textbf{Back to the full CMM (0D):}
      \begin{itemize}\itemsep3pt
        \item spring $\leftrightarrow$ constituent law $\sigma^k(\lambda_e^k)$
        \item dashpot $\leftrightarrow$ mass turnover / removal $k_d^k$
        \item motor $\leftrightarrow$ deposition stretch $G^k$
      \end{itemize}
    \end{column}
  \end{columns}
  \srccite{Cyron, Aydin, Humphrey, Biomech Model Mechanobiol (2016)}
\end{frame}

\begin{frame}[t]{Growth (Mass) and Remodeling (Natural Configuration)}
  \textbf{Growth}: mass production $\propto$ current mass, driven by stress
  deviation. \textbf{Remodeling}: the mean natural configuration evolves as new mass
  enters at $\lambda/G^k$ and old mass leaves at the current value.
  \dlab{general form --- mass \& mean natural-config \emph{tensor}}
  \vspace{-0.2em}
  \giveneq{\dot M^k = M^k\,(K_\sigma^k k_d^k)\!\left(\tfrac{\bar\sigma}{\bar\sigma_h}-1\right),
    \qquad \dot{\bm{F}}_n^k = f\big(\bm{F},\bm{F}_n^k;\,m^k,M^k\big).}
  \zlab{reduction --- scalar mass and scalar natural stretch}
  \vspace{-0.2em}
  \begin{columns}[c]
    \begin{column}{0.5\textwidth}
      \slboxeq{2.2em}{eq:homog_mass}{\frac{\dd M^k}{\dd t}
        = M^k K_\sigma^k k_d^k\!\left(\tfrac{\bar\sigma}{\bar\sigma_h}-1\right)}
    \end{column}
    \begin{column}{0.5\textwidth}
      \slboxeq{2.2em}{eq:homog_remodel}{\frac{\dd\lambda_n^k}{\dd t}
        = \frac{m^k}{M^k}\!\left(\frac{\lambda}{G^k}-\lambda_n^k\right)}
    \end{column}
  \end{columns}
  \vspace{0.2em}
  Holding $\lambda$ fixed, remodeling relaxes constituent stress exponentially to
  $\sigma_h^k$ ($\sim 1/k_d$) --- not to zero, but to a prestress.
  \srccite{Cyron, Aydin, Humphrey, Biomech Model Mechanobiol (2016)}
\end{frame}

\begin{frame}[t]{It Tracks the Full Model Closely --- Much Cheaper}
  Two ODEs per constituent replace the full deposition history. The homogenized
  trajectory sits almost on the full constrained mixture, while kinematic growth
  under-responds to elastin loss --- it grows constituents in lockstep and cannot
  remodel their natural configurations individually.
  \vfill
  \begin{center}\color{AccentRed}
  The homogenized model follows the full model's \emph{time trajectory} closely,
  at $O(N)$ instead of $O(N^2)$.
  \end{center}
\end{frame}

% ---- pen & paper: when are the two models close? ---------------------------
\begin{frame}[t]{}
  \begin{yourturnbox}{When do homogenized and full agree?}
    \qq{The homogenized model keeps only the \emph{mean} natural configuration per
    constituent. Under what assumptions should it match the full deposition-history
    model --- and when would you expect them to diverge? \qkind{discuss}}
    \qres{The two remodeling pictures: full heredity integral vs.\ the single mean
    stretch $\lambda_n^k$ \eqref{eq:homog_remodel}.}
    \qhint{Think about the \emph{spread} of natural configurations across cohorts.
    When is a single mean a good summary of the whole history?}
  \end{yourturnbox}
\end{frame}

\begin{frame}[t]{}
  \begin{answerbox}{When do homogenized and full agree?}
    They agree when each constituent's cohorts share \emph{nearly the same} natural
    configuration --- i.e.\ slow, gradual loading where deposition stretch $G^k$ and
    turnover are roughly stationary, so the cohort spread stays narrow and the mean
    represents it well. {\color{ABlue}\bfseries They diverge} under fast or strongly
    time-varying insults, when cohorts deposited at very different states coexist and
    a single mean stretch can no longer capture the distribution.
  \end{answerbox}
\end{frame}

% ---- example + run: ex04 ---------------------------------------------------
\begin{frame}[t]{}
  \begin{setupbox}{Example --- homogenized vs.\ full (\texttt{ex04})}
    \sprob{Run both mixture theories on the \emph{same} artery and insult. The
    homogenized model keeps one mean natural stretch per constituent (two ODEs,
    \eqref{eq:homog_mass}--\eqref{eq:homog_remodel}) instead of the full cohort
    history. How closely does it track the full model, and how much faster is it?}
    \srun{\scodeline{uv run python run.py configs/ex04\_homogenized\_vs\_full.yaml}
      Here \texttt{study: compare} runs
      \texttt{theories: [constrained\_mixture, homogenized\_cmm]} on one scenario
      and prints each runtime and the largest mass-ratio gap between them.}
    \sstudy{Stress the comparison with a harsher, longer insult:
      \scodeline{... --set insult.elastin\_surviving=0.3 --set insult.ramp=10}
      \scodeline{... --set simulate.t\_end=4000\ \ \ \ \ \# does the agreement hold?}}
  \end{setupbox}
\end{frame}

% ---- interactive coding exercise: ex04 -------------------------------------
\begin{frame}[t]{}
  \begin{yourturnbox}{Homogenized vs.\ full}
    \qq{Is the difference between the two curves visible by eye? By what factor is
    the homogenized model faster --- and does the agreement survive a harsher,
    longer insult?}
    \qhint{Read the printed mass-ratio gap and the two runtimes. The homogenized
    model should track the full one closely wherever the loading is slow and
    gradual.}
  \end{yourturnbox}
\end{frame}

\begin{frame}[t]{}
  \begin{answerbox}{Homogenized vs.\ full}
    The homogenized curve overlays the full CMM closely at a fraction of the
    runtime ($O(N)$ vs.\ $O(N^2)$). {\color{ABlue}\bfseries Learning:} it reproduces
    a \textbf{similar time trajectory} to the full model --- cheaply.
  \end{answerbox}
  \vspace{-0.4em}
  \bigfigw{fig04_compare_all.pdf}{0.52\linewidth}{Your overlay: homogenized (dashed)
  on the full CMM (black); kinematic growth (blue) under-responds.}
\end{frame}
